\documentclass[10pt,a4paper]{amsart}
\usepackage[margin=19mm]{geometry}
\usepackage{amsmath,amssymb,mathtools}
\usepackage[scr=rsfs,cal=euler]{mathalfa}
\usepackage{xfrac}
\usepackage[unicode,hidelinks,bookmarks=false]{hyperref}
\newcommand{\ie}{i.e.\ }
\newcommand{\cf}{cf.\ }
\newcommand{\resp}{respectively\ }
\newcommand{\Subsection}{\subsection}
\providecommand{\nobreakdash}{\nobreak\hbox{-}}
\setlength{\emergencystretch}{2em}
\multlinegap0pt
\usepackage{bm}
\usepackage{bbm}
\usepackage{dsfont}
\usepackage{xspace}
\usepackage{cancel}
\usepackage{mathtools}
\usepackage{amsthm}
\makeatletter
\@ifundefined{Theorem}{\newtheorem{Theorem}{Theorem}}{}
\@ifundefined{Lemma}{\newtheorem{Lemma}[Theorem]{Lemma}}{}
\makeatother
\makeatletter
\newcommand{\abs}[1]{\lvert #1 \rvert}          % Formatting for the absolute value
\newcommand{\eps}{\varepsilon}
\newcommand{\ol}[1]{\overline{#1}}
\newcommand{\supp}{\mathrm{supp}}
\expandafter\ifx\csname tehtratk\endcsname\relax
\newenvironment{tehtratk}%
             {\begin{list}{\arabic{enumi}.}{\usecounter{enumi}%
              \setlength{\labelsep}{0.5em}%
              \settowidth{\labelwidth}{(\arabic{enumi})}%
              \setlength{\leftmargin}{\labelwidth+\labelsep}}}%
             {\end{list}}
\fi
\makeatother
\title[A free boundary approach to non-scattering obstacles with va]{A free boundary approach to non-scattering obstacles with vanishing contrast}
\author{Mikko Salo, Henrik Shahgholian}
\date{Source: arXiv:2506.22328v1 (CC BY 4.0)}
\begin{document}
\maketitle

\noindent\emph{Source and licence.} A free boundary approach to non-scattering obstacles with vanishing contrast. Version arXiv:2506.22328v1 (CC BY 4.0), distributed under the Creative Commons Attribution 4.0 International licence, as recorded in the arXiv licence metadata for this exact version and shown on the abstract page https://arxiv.org/abs/2506.22328v1. Full rights notice: licenses/arxiv-2506.22328-CC-BY-4.0.txt.\par\smallskip

\noindent\textit{Editorial scope.} Counted unit: the Lemma whose source label is \texttt{lemma\_nondegeneracy} in the article (source0.tex lines 1071--1076, source label \texttt{lemma\_nondegeneracy}), delivered together with the article's own complete original proof of it (source0.tex lines 1077--1123, one \texttt{proof} environment of 47 lines). No mathematics is written, repaired or re-derived: statement and proof are the article's own text line for line. Required context retained uncounted: none. Every definition, display and label of those lines is retained unchanged. Omitted from this package and not redistributed: 1--1070, 1124--1911. Nothing inside a retained excerpt is omitted or altered: every retained body is delivered exactly as the article's lines are written, so no omission receipt of any kind accompanies this package. Citations: the retained excerpts carry no citation command at all, so no bibliography entry of the article is transcribed and no reference is redistributed. Reference adaptation: none was needed and none was made; every reference inside the delivered unit resolves inside it, so no unresolved reference or citation is produced. Printed slips kept literal: no printed word, symbol, hypothesis or number of the counted statement or of its proof was altered. Layout and class: the article's own class is \texttt{amsart} with packages including amsmath, amscd, amssymb, amsthm, amsrefs, hyperref, graphicx, mathrsfs, calc, comment, xcolor; the wrapper is the lane's pinned \texttt{amsart} editorial wrapper at 10pt on A4, which restates the theorem-like environments the retained text needs and adds the article's own macro definitions that the retained text needs (\texttt{\textbackslash abs}, \texttt{\textbackslash eps}, \texttt{\textbackslash ol}, \texttt{\textbackslash supp}), with extra wrapper packages (bm, bbm, dsfont, xspace, cancel, mathtools). The delivered numbering is the unit's own. Assets: no retained span contains a figure environment, a graphics-inclusion command, a PDF-page inclusion, a table or a TikZ picture; no image file is copied into this package, the package declares no asset dependency and no image file is redistributed with it. Licence: version arXiv:2506.22328v1 of this article is distributed under the Creative Commons Attribution 4.0 International licence, as recorded in the arXiv licence metadata for this exact version and shown on the abstract page https://arxiv.org/abs/2506.22328v1. A full rights notice is delivered at licenses/arxiv-2506.22328-CC-BY-4.0.txt. Characters: every retained excerpt is pure ASCII, so the delivered engine, XeLaTeX with its default Latin Modern text font, reports no missing character.
\begin{Lemma} \label{lemma_nondegeneracy}
For any $\eps$ with $0 < \eps < 1$ there is $c_{\eps} > 0$ so that 
\[
\sup_{B(x, \eps \abs{x})} \abs{u} \geq c_{\eps} \abs{x}^{m+2} \text{ whenever $x \in \ol{D} \cap \ol{B}_{r_0}$}.
\]
\end{Lemma}

\begin{proof}

The condition for $\supp(f)$ ensures that $\ol{B}_{r_0} \cap \ol{D} \subset \supp(\Delta u)$ and hence also $\ol{B}_{r_0} \cap \ol{D} \subset \supp(u)$. In particular, $\supp(u) \cap \ol{B}_{r_0} = \ol{D} \cap \ol{B}_{r_0}$. Since $D$ is a Lipschitz domain, there is $\delta > 0$ so that possibly after shrinking $r_0$ we have the following consequence of the interior cone property:
\begin{equation} \label{interior_ball}
\text{If $x \in \ol{D} \cap \ol{B}_{r_0}$, then $B(x, r) \cap \ol{D}$ contains a ball of radius $\delta r$ whenever $r \leq r_0$.}
\end{equation}

We now prove the lemma. If the claim were not true, then there would be some $\eps \in (0,1)$ so that for any $j \geq 1$ there is $x_j \in \ol{D} \cap \ol{B}_{r_0}$, $x_j \neq 0$, so that 
\begin{equation} \label{nondegeneracy_contradiction}
\sup_{B(x_j, \eps \abs{x_j})} \abs{u} < \frac{1}{j} \abs{x_j}^{m+2}.
\end{equation}
If one had $\abs{x_j} \geq c > 0$ for all $j$, then after passing to a subsequence there would be $x \in \ol{D} \cap \ol{B}_{r_0}$ with $\abs{x} \geq c$ so that $u|_{B(x, \eps\abs{x}/2)} = 0$. This would contradict the fact that $\supp(u) \cap \ol{B}_{r_0} = \ol{D} \cap \ol{B}_{r_0}$. Thus we may assume, after taking a subsequence, that $x_j \to 0$.

Let $r_j = \abs{x_j}$ and define the rescaled function 
\[
u_{r_j}(x) = \frac{u(r_j x)}{r_j^{m+2}}.
\]
Now using \eqref{interior_ball} with $x=x_j$ and $r = \eps r_j$, for any $j \geq 1$ there is $y_j$ so that 
\begin{equation} \label{nondegeneracy_interior_ball_second}
B(y_j, \delta \eps r_j) \subset B(x_j, \eps r_j) \cap \ol{D}.
\end{equation}
It also follows that $(1-\eps) r_j \leq \abs{y_j} \leq (1+\eps) r_j$. Then by \eqref{nondegeneracy_contradiction} we have 
\[
\sup_{B(y_j, \delta \eps r_j)} \abs{u} < \frac{1}{j} r_j^{m+2}.
\]
Writing $z_j = y_j/r_j$, this means that 
\[
\sup_{B(z_j, \delta \eps)} \abs{u_{r_j}} < \frac{1}{j}.
\]
By taking another subsequence we may assume that $z_j$ converges to some $z$ with $1-\eps \leq \abs{z} \leq 1+\eps$. Thus we have 
\begin{equation} \label{nondegeneracy_uj_convergence}
u_{r_j} \to 0 \text{ in $L^{\infty}(B(z, \delta \eps/2))$}.
\end{equation}
On the other hand we have 
\[
\Delta u_{r_j}(x) = r_j^{-m} \Delta u(r_j x) = -q r_j^2 u_{r_j}(x) + (P(x) + r_j^{-m} R(r_j x)) \chi_D(r_j x).
\]
Now \eqref{nondegeneracy_interior_ball_second} implies that 
\[
\Delta u_{r_j}(x) = -q r_j^2 u_{r_j}(x) + P(x) + r_j^{-m} R(r_j x) \text{ in $B(z_j, \delta \eps)$}.
\]
Since $r_j \to 0$ and $\abs{u_{r_j}} \leq C$, $\abs{R} \leq C \abs{x}^{m+\alpha}$, we have 
\[
\Delta u_{r_j} \to P \text{ in $L^{\infty}(B(z, \delta \eps/2))$}.
\]
Since $P \not\equiv 0$, this contradicts \eqref{nondegeneracy_uj_convergence} by the uniqueness of distributional limits.
\end{proof}

\end{document}
